\section{Pre-Registered Predictions}
\label{sec:predictions}

\subsection{Mathematical Derivation}

The joint state space of the non-ergodic mixture is $(k, s) \in \{1,\ldots,K\} \times \{1,2,3\}$, with $3K$ total states. Since component identity $k$ is fixed within each sequence, the predictive vector factorizes:
\begin{equation}
    P(k, s \mid z_{1:t}) = P(k \mid z_{1:t}) \cdot P(s \mid k, z_{1:t})
\end{equation}
where $P(k \mid z_{1:t})$ is the \emph{meta-belief} (posterior over components) and $P(s \mid k, z_{1:t})$ is the \emph{within-component belief state}.

The joint transition operator is block-diagonal:
\begin{equation}
    T^{(z)} = \bigoplus_{k=1}^K P(k) \cdot T_k^{(z)}
\end{equation}
This is \emph{not} a tensor product $\bigotimes_n T_n^{(z)}$, so the standard conditionally-independent factored representation from~\cite{shai2026factored} does not directly apply. However, the structure admits a natural factorization because $k$ is static: after sufficient context, $P(k \mid z_{1:t})$ concentrates on the true component, and the predictive vector effectively reduces to the single-component case.

\paragraph{Meta-belief update.} The posterior over components evolves as:
\begin{equation}
    P(k \mid z_{1:t}) \propto P(k) \prod_{\tau=1}^{t} P(z_\tau \mid z_{1:\tau-1}, k)
\end{equation}
where $P(z_\tau \mid z_{1:\tau-1}, k) = \veta_k^{(z_{1:\tau-1})} T_k^{(z_\tau)} \mathbf{1}$ is the predictive probability under component $k$. The rate at which the meta-belief concentrates depends on the KL divergence between component predictive distributions---larger divergences enable faster source discrimination.

\subsection{Geometric Predictions}

\paragraph{Prediction 1: Dimensionality.}
We predict the effective dimensionality of the residual stream representation should be approximately $K + 1$ under a factored representation:
\begin{itemize}
    \item $(K-1)$ dimensions for the meta-belief simplex over $K$ components
    \item $2$ dimensions for within-component belief (the 2-simplex shared across components)
\end{itemize}
For $K = 3$: factored prediction is $4$ dimensions. The joint alternative requires $3K - 1 = 8$ dimensions. Separated gaskets would use $\sim 2K = 6$ dimensions.

\paragraph{Prediction 2: Context-position dependence.}
At position $t = 1$, the meta-belief is near-uniform: all components are equally plausible. The geometry should resemble a superposition of $K$ overlapping gaskets. As $t$ increases, the meta-belief sharpens, and the geometry should resolve into component-specific clusters, each containing its own gasket structure. We expect the per-position loss to decrease with context, approaching the Bayes-optimal rate, with the rate of decrease governed by the KL divergence between components.

\paragraph{Prediction 3: Layer dependence.}
Based on the layer-by-layer belief construction observed in~\cite{shai2024transformers}:
\begin{itemize}
    \item \textbf{Embedding}: 3 points (one per token), no component or belief structure.
    \item \textbf{Early layers}: Component discrimination should emerge first---the model identifies which source is active by comparing token statistics against learned priors.
    \item \textbf{Later layers}: Within-component gasket structure should sharpen as the model refines its belief state conditioned on the identified component.
\end{itemize}
This ordering reflects the computational dependency: you must know \emph{which} process to use the \emph{right} transition matrices for belief updating.

\subsection{Geometric Intuition}

Multiple geometries are conceivable:
\begin{enumerate}
    \item \textbf{Separated gaskets}: $K$ distinct Sierpinski gaskets in different regions of activation space, with large offsets encoding component identity. Total: $\sim 2K$ dimensions.
    \item \textbf{Factored (shared belief subspace)}: Component identity in a $(K-1)$-dimensional subspace, with all gaskets superimposed in a shared 2D belief subspace. Total: $K+1$ dimensions. This is the strongest form of factoring.
    \item \textbf{Conditional (component-specific belief subspaces)}: Each component gets its own 2D belief subspace with different orientations, plus a shared component-identity subspace. Total: $2K + (K-1)$ dimensions.
\end{enumerate}

Our formal prediction favors geometries 1 or 2. Geometry 1 is natural given the block-diagonal (not tensor-product) structure. Geometry 2 would be evidence that the factoring inductive bias extends beyond the conditionally-independent case to block-diagonal generators.

We will not be surprised if the model initially learns geometry 1 and transitions toward geometry 2 over training, following the inductive bias toward dimensional efficiency observed in~\cite{shai2026factored}.
